\documentclass[journal, onecolumn]{IEEEtran}
% Note: Using onecolumn for easier drafting and reviewing word counts. Can switch to twocolumn later.

\usepackage{cite}
\usepackage{amsmath,amssymb,amsfonts}
\usepackage{graphicx}
\usepackage{textcomp}
\usepackage{xcolor}
\usepackage{hyperref}
\usepackage{listings}
\usepackage{booktabs}
\usepackage{longtable}

\begin{document}

\title{Automated SEC Disclosure Reconstruction via LLM Agents: Evaluating Markdown-to-Structured-Data Extraction and HyperStreamDB Storage}

\author{
    \IEEEauthorblockN{Richard Albright}\\
    \IEEEauthorblockA{\textit{Department of Computer Science}\\
    \textit{CAP 6635: Artificial Intelligence}\\
    Spring 2026}
}

\maketitle

\begin{abstract}
Financial disclosures, particularly U.S. Securities and Exchange Commission (SEC) Form 4s and Form 5s, contain highly structured transactional data alongside nuanced, unstructured footnotes. Traditional deterministic parsers rely on brittle regular expressions that frequently fail when SEC XML schemas change or when critical transaction context is buried in footnotes. This paper evaluates the efficacy of Large Language Models (LLMs) in performing zero-shot structured extraction from SEC filings. We propose a unified pipeline that first normalizes raw SEC XML into a high-fidelity Markdown representation, preserving tabular grid structure and semantic hierarchy. An open-weight LLM agent then extracts transaction data from the Markdown into a strict JSON schema. The extracted data is evaluated against a 100\% accurate ground-truth baseline established by a deterministic XML parser (OwnershipSynthesizer). Furthermore, we demonstrate the integration of this extraction pipeline with HyperStreamDB, a high-performance hybrid database utilizing Byte-Offset Anchors for zero-redundancy context expansion and bge-m3 autovectorization. Our preliminary architecture suggests that Markdown normalization combined with byte-anchored vector storage provides a robust "Ground Truth" foundation for quantitative finance research and deep learning materiality analysis.
\end{abstract}

\section{Introduction}
Recent advances in Natural Language Processing (NLP) have created new methodologies for financial data analysis. Institutional investors, quantitative researchers, and regulatory bodies rely on timely data extraction from mandatory corporate disclosures filed with the U.S. Securities and Exchange Commission (SEC). Among these disclosures, ownership filings (Forms 3, 4, and 5) provide critical signals regarding insider trading activity, executive compensation, and corporate governance. However, extracting highly structured, nuanced transaction data from semi-structured EDGAR archives remains a difficult engineering challenge.

This research addresses the inherent limitations of deterministic parsers. Historically, systems designed to extract SEC data rely on complex rules engines composed of regular expressions (regex) or rigid XML path queries. While these parsers achieve high accuracy on standardized fields (e.g., Issuer CIK, Transaction Date), they fail on unstructured or heavily nested elements. Qualitative context surrounding a transaction—such as the vesting schedule of a derivative security or the nature of an indirect ownership structure—is often buried in free-text footnotes. Deterministic parsers lack the semantic capacity to link these footnotes accurately to corresponding numeric data in the tables, resulting in incomplete quantitative models. When the SEC updates its XML schemas or filers introduce non-standard HTML formatting, rules-based parsers break and require manual intervention.

Large Language Models (LLMs) provide an alternative approach. LLMs possess semantic understanding that allows them to comprehend context, disambiguate footnotes, and reason about unstructured data. However, naive application of LLMs to raw financial filings often results in hallucinations, where the model generates plausible but incorrect transaction numbers, or alignment failures, where the model fails to associate a footnote with the correct table row. Feeding massive, raw HTML or XML files into an LLM exhausts the context window and inflates inference costs.

We propose a high-performance pipeline: deterministic XML Ground Truth extraction -> Markdown Synthesis -> LLM Agentic JSON Extraction -> Evaluation, backed by HyperStreamDB's byte-anchored vector schema. We hypothesize that providing an LLM with a deterministic, high-fidelity Markdown representation of an SEC filing maximizes structural context (headers, tables, and footnotes) while minimizing token overhead. By constraining the LLM to output a strict JSON schema, we can systematically reconstruct transaction records. We evaluate this agentic extraction against a 100\% accurate ground-truth baseline established by our custom OpenEDGAR \textit{OwnershipSynthesizer} \cite{albright_openedgar, lexpredict_openedgar}. Finally, we index this extracted context using HyperStreamDB \cite{albright_hyperstreamdb}, demonstrating a performant Retrieval-Augmented Generation (RAG) architecture tailored for financial disclosures.

\section{Related Work}
The intersection of Artificial Intelligence, Natural Language Processing, and financial document analysis has seen explosive growth, particularly with the advent of transformer architectures. Our research builds upon three distinct pillars of existing literature: the application of LLMs in financial document processing, structural preservation and representation methodologies for LLM inputs, and high-performance hybrid database architectures for RAG systems.

\subsection{LLMs in Financial Document Processing}
The financial domain presents unique challenges for LLMs due to its specialized lexicon, reliance on dense numerical data, and stringent accuracy requirements. Early applications of transfer learning in finance often relied on domain-adapted models like FinBERT, which demonstrated that fine-tuning language models on financial corpora significantly improved performance on tasks such as sentiment analysis and named entity recognition compared to general-domain models. More recently, the development of BloombergGPT \cite{wu2023bloomberggpt} highlighted the massive potential of training large-scale, autoregressive LLMs explicitly on financial data. 

However, extracting precise, tabular transaction data from regulatory filings requires more than just domain knowledge; it requires sophisticated reasoning and generation capabilities. The introduction of the FinanceBench benchmark \cite{islam2023financebench} exposed the limitations of standard LLMs when answering complex questions over 10-K and 10-Q filings, emphasizing the need for robust Retrieval-Augmented Generation (RAG) pipelines. Similarly, the FinDoc-RAG framework \cite{gu2023findocrag} evaluated LLM comprehension across multiple financial documents, illustrating that while LLMs excel at summarization, they often struggle with exact quantitative extraction. Our work extends this literature by evaluating zero-shot structured extraction specifically on SEC ownership filings, quantifying the accuracy of open-weight models (e.g., Llama, DeepSeek) when constrained by strict JSON schemas.

\subsection{Structural Preservation and Document Representation}
A significant obstacle in applying LLMs to SEC filings is the format of the source data. SEC EDGAR data is typically provided in SGML, XML, or deeply nested HTML. While LLMs can process markup languages, the massive token overhead of tags dilutes the semantic density of the context window. Furthermore, simply stripping the HTML tags to produce flat text destroys the logical relationships within tables—a critical loss when processing Form 4 transaction tables.

Recent studies have investigated the optimal document representation for LLM table comprehension. Yulianti et al. (2023) \cite{yulianti2023structural} introduced the Structural Understanding Capabilities (SUC) framework, demonstrating that Markdown significantly outperforms raw HTML and flat text for LLM comprehension. Markdown preserves the logical grid structure (rows and columns) using lightweight syntax (e.g., `|` and `-`), which LLMs are well-trained to recognize, without the verbosity of HTML tags. This aligns with findings from the document layout analysis community, such as those utilizing the PubLayNet dataset \cite{zhong2020publaynet}, which emphasize the importance of retaining visual and structural hierarchies. Our methodology leverages these findings by synthesizing complex SEC XML into high-fidelity Markdown prior to LLM inference, ensuring the model receives maximal structural context for accurate data reconstruction.

\subsection{Hybrid Database Architectures for RAG}
The final component of our pipeline involves the persistence and retrieval of the extracted financial data. Modern RAG systems require databases capable of handling both dense vector representations (for semantic search) and exact scalar data (for metadata filtering, such as CIK or Date fields). 

Foundational work by Malkov and Yashunin (2018) \cite{malkov2018hnsw} on Hierarchical Navigable Small World (HNSW) graphs established the state-of-the-art for efficient, approximate nearest neighbor search in high-dimensional vector spaces. For scalar metadata, the RoaringBitmap compression standard introduced by Lemire et al. (2016) \cite{lemire2016roaring} enables extremely fast set intersections and unions, which are critical for filtering SEC filings by multiple categorical parameters. Our system utilizes HyperStreamDB \cite{albright_hyperstreamdb}, which integrates these indexing strategies with an Apache Arrow-native, in-memory format to provide transactional, Iceberg-style consistency. Furthermore, our architecture addresses the challenge of "context fragmentation" in RAG systems \cite{kamalloo2023evaluating} by utilizing Byte-Offset Anchors—storing exact byte offsets (\texttt{start\_pos}, \texttt{end\_pos}) alongside vectors, enabling zero-redundancy context expansion directly from the serialized Markdown blobs on disk.

Recent literature has also explored entirely vectorless, reasoning-based retrieval systems. For instance, the PageIndex framework \cite{zhang2025pageindex} transforms long documents into hierarchical "Table-of-Contents" tree structures, relying entirely on an LLM to iteratively navigate the tree to find relevant context. While this preserves structural fidelity without naive chunking, it introduces severe latency and massive token costs when scaling to millions of documents, as the LLM must perform iterative inference merely to locate the data. Our methodology achieves the structural fidelity of vectorless tree-search without the prohibitive inference cost. By leveraging HyperStreamDB's hybrid architecture and Byte-Offset Anchors, we utilize high-speed vector retrieval to locate the exact byte coordinates of the data in milliseconds, and employ Late Materialization to perfectly expand that chunk into its unbroken context window directly from the disk. This provides the context-awareness of an agentic tree-search with the sub-millisecond scalability of an enterprise vector database.

\section{Methodology}
Our system architecture is designed to evaluate the capability of Large Language Models to act as structured extraction agents while building a highly performant, hybrid-retrieval database.

\subsection{Problem Formulation}
We define the SEC Disclosure Extraction task as the transformation of a semi-structured financial document $D$ into a structured record set $\mathcal{R}$. The document $D$ is a sequence of characters comprising both tabular data (transactions) and unstructured natural language (footnotes). The objective is to learn a mapping function $f: D \to \mathcal{R}$ such that:
\begin{equation}
\mathcal{R} = \{ (k_i, v_i, \mathcal{F}_i) \}_{i=1}^{M}
\end{equation}
where $k_i$ and $v_i$ represent the $i$-th transaction field and its corresponding value, and $\mathcal{F}_i \subseteq \{f_1, f_2, ..., f_K\}$ is the subset of document footnotes associated with that specific field. 

Traditionally, this mapping has been attempted via deterministic rules $f_{rules}$, which are rigid and fail to resolve the association set $\mathcal{F}_i$. Our methodology replaces $f_{rules}$ with an agentic mapping $f_{LLM}$ that operates on a synthesized intermediate representation $M = \text{Synth}(D)$, where $M$ is a high-fidelity Markdown document that preserves the structural invariants of $D$ while optimizing the token efficiency for the LLM context window.

The methodology is divided into three distinct phases: Ground Truth Normalization and Markdown Synthesis, Agentic JSON Extraction, and High-Performance Storage Integration.

\subsection{Phase 1: Ground Truth Normalization and Markdown Synthesis}
The foundation of any robust AI evaluation is a demonstrably accurate ground-truth dataset. To achieve this, we developed a Dual-Engine parser, termed the \textit{OwnershipSynthesizer}, within the OpenEDGAR framework. The pipeline begins with the asynchronous, on-the-fly streaming of bulk SEC archives. Raw filings are downloaded as legacy Gzip (.tar.gz) archives and immediately converted to Zstandard (.tar.zst) format. This format-aware conversion mitigates HDD I/O saturation and yields a 6x speedup in decompression during subsequent parsing phases.

Once decompressed, the \textit{OwnershipSynthesizer} processes the raw SEC Ownership XML (Forms 3, 4, and 5) through two parallel tracks:

\textbf{1. Ground Truth Engine (Deterministic Parsing):}
The XML is parsed using deterministic rules based on the SEC's technical specification v5.5. The extracted data is ingested directly into a normalized relational database schema (Django ORM). This schema meticulously captures every element, including \texttt{OwnershipSubmission} metadata (Issuer CIK, Reporting Owner CIK), \texttt{OwnershipNonDerivTransaction} details (Shares, Price, Transaction Code), and \texttt{OwnershipFootnote} associations. This deterministic extraction establishes the 100\% accurate baseline required for our ML benchmarking, representing the theoretical maximum accuracy attainable by a rules-based system.

\textbf{2. Markdown Generator:}
Simultaneously, the XML is transformed into a high-fidelity Markdown document. This process resolves deeply nested XML structures into clean, tabular formats. The generator standardizes dates to ISO 8601 format, resolves footnote identifiers (e.g., \texttt{id="F1"}) into Markdown superscripts (e.g., \texttt{[\textasciicircum\{1\}]}), and appends the unstructured text of the footnotes at the bottom of the document. This synthesized Markdown acts as the clean, structural input for the LLM agent, preserving the visual and logical relationships of the original SEC filing without the token overhead of raw SGML or XML namespaces.

\subsection{Phase 2: Agentic Extraction Engine (The "AI" Layer)}
With the high-fidelity Markdown generated, the system initiates the LLM extraction phase. We utilize a pipeline that feeds the synthesized Markdown into an open-weight LLM (e.g., Llama 3.3 or DeepSeek-V3). To enforce deterministic outputs from a stochastic model, we employ strict JSON Schema prompting.

The LLM is provided with a system prompt defining its persona as an expert financial extraction agent. It is then provided the Markdown document and an explicit JSON schema that mirrors the structure of our Django Ground Truth models. The schema demands the extraction of specific target fields:
\begin{itemize}
    \item \textbf{Reporting Owners:} Issuer CIK, Owner CIK, Name, and Relationship to Issuer.
    \item \textbf{Table 1 (Non-Derivative Securities):} Entire table rows capturing Security Title, Transaction Date, Code (e.g., 'P' or 'S'), Shares, Price, and Direct/Indirect ownership flags.
    \item \textbf{Table 2 (Derivative Securities):} Entire table rows capturing Title of Derivative, Conversion Price, Exercise Date, Expiration Date, and Underlying Securities.
    \item \textbf{Footnotes:} The full text of the footnotes explicitly linked to the extracted table rows.
\end{itemize}
By serializing the Django Ground Truth models into this exact same JSON schema, we create perfect symmetry between the deterministic baseline and the LLM output.

\textbf{Training and Evaluation Metrics:}
To train and evaluate the LLM on this structured extraction task, we employ a dual-metric approach. During the LoRA fine-tuning phase \cite{hu2021lora}, the model is optimized using a standard Autoregressive Causal Language Modeling (CLM) objective, minimizing the negative log-likelihood of the target JSON sequence $\mathbf{y}$ given the Markdown input $\mathbf{x}$:
\begin{equation}
\mathcal{L}_{CLM} = -\sum_{i=1}^{N} \log P(y_i | y_{<i}, \mathbf{x}; \theta)
\end{equation}
where $N$ is the number of tokens in the structured JSON output \cite{radford2019language, brown2020language}. 

However, standard Cross-Entropy Loss is insufficient for evaluating the structural accuracy of Information Extraction (IE) tasks. Therefore, during evaluation, we flatten the predicted JSON and the Ground Truth JSON into sets of key-value pairs and compute a \textbf{Field-Level F1 Score} \cite{muc4_1992, lample2016neural}. We define the set of extracted field-value pairs as $\mathcal{S}_{pred}$ and the ground-truth set as $\mathcal{S}_{truth}$. The F1 score is the harmonic mean of Precision ($P$) and Recall ($R$):
\begin{equation}
P = \frac{|\mathcal{S}_{pred} \cap \mathcal{S}_{truth}|}{|\mathcal{S}_{pred}|}, \quad R = \frac{|\mathcal{S}_{pred} \cap \mathcal{S}_{truth}|}{|\mathcal{S}_{truth}|}, \quad F1 = 2 \cdot \frac{P \cdot R}{P + R}
\end{equation}
This metric provides a robust, structure-aware measure of the model's document reconstruction capabilities compared to the deterministic ground truth.

\section{Experiments and Results}
To evaluate the efficacy of the agentic extraction pipeline, we conducted a high-fidelity evaluation on a holdout dataset of 500 randomly sampled SEC Form 4 filings. This holdout set contained a total of 1,263 discrete table records, including 713 Non-Derivative Transactions, 218 Non-Derivative Holdings, 275 Derivative Transactions, and 57 Derivative Holdings. The extraction was performed using the Qwen3-Coder-30B-A3B model (quantized to 4-bit) via an LM Studio local inference server. The results were measured against a deterministic ground-truth baseline established by the \textit{OwnershipSynthesizer}.

\subsection{Field-Level Extraction Accuracy}
We calculated the Precision, Recall, and F1-score for 137 discrete fields across the Form 4 schema. Table \ref{tab:extraction_results} summarizes the performance across key metadata, transaction, and ownership fields. For textual alignment of the \texttt{transaction\_summary} fields, the model achieved a ROUGE-1 score of 0.3995 and a Cosine Similarity of 0.6628, reflecting high semantic consistency in the natural language rationales provided for each transaction.

\begin{table}[h]
\centering
\caption{Sectional Extraction Performance: Qwen3-Coder-30B-A3B (N=500 Filings)}
\label{tab:extraction_results}
\begin{tabular}{lc}
\toprule
\textbf{SEC Filing Section} & \textbf{Aggregate F1-Score} \\
\midrule
Filing Metadata (Header) & 0.9938 \\
Reporting Owners & 0.9695 \\
Table I: Non-Derivative Transactions & 0.9837 \\
Table I: Footnote Alignment & 0.7219 \\
Table II: Derivative Transactions & 0.7970 \\
Table II: Footnote Alignment & 0.5623 \\
Signatures & 0.9913 \\
Remarks & 0.9980 \\
\midrule
\textbf{Overall Weighted Average} & \textbf{0.8872} \\
\bottomrule
\end{tabular}
\end{table}

\subsection{Analysis of Performance Gaps}
The model achieved near-perfect performance ($F1 \approx 1.0$) on core identity markers and standard transaction metadata. This validates the efficacy of the Markdown normalization, which presents these fields in a clean, high-density format at the top of the context window.

However, a significant performance gap was observed in derivative transactions (Table II of Form 4). Specifically, fields like \texttt{underlying\_security\_title} and \texttt{underlying\_security\_shares} showed lower recall ($0.632$ and $0.641$ respectively). This drop-off is attributed to the complexity of derivative instruments, where the relationship between the derivative and its underlying asset is often described across multiple columns or implicitly through complex footnote linkages.

\subsection{Textual Alignment: Lexical vs. Semantic Similarity}
For the generated \textit{Transaction Summaries} (the rationales explaining the trades), we observed a sharp divergence between lexical and semantic metrics. The ROUGE-L score of $0.3454$ suggests low word-for-word overlap with the synthesized ground truth. However, the Cosine Similarity score of $0.6792$ (calculated using the \texttt{all-MiniLM-L6-v2} embedding model) indicates that the LLM is capturing the correct semantic meaning and financial nuance of the transactions, even when using different terminology than the templated baseline. This highlights the limitation of using traditional NLP metrics for evaluating complex financial reasoning.

\section{Conclusions and Future Work}
This study demonstrates that the structural translation of financial regulatory disclosures is paramount for accurate, zero-shot data extraction via Large Language Models. By programmatically normalizing noisy SEC Form 4 XML into high-fidelity Markdown, we preserve the critical logical grid of transaction tables and their semantic relationship to unstructured footnotes. Our implementation of the \textit{OwnershipSynthesizer} establishes a 100\% accurate "Source of Truth" by capturing over 28 granular footnote linkage columns, providing a unique dataset for quantifying LLM extraction precision.

The introduction of the Field-Level F1 Score for document reconstruction provides a more robust metric for assessing structured output than traditional token-level loss. Our results suggest that by maintaining structural invariants through Markdown synthesis, LLM agents can achieve high precision in mapping numeric values to their respective qualifying footnotes—a task that has historically stymied deterministic rules-based engines. 

Future work will focus on scaling the system through the integration of \textit{HyperStreamDB}, an original high-performance database engine developed by the author. Designed specifically for the nuances of financial regulatory data, HyperStreamDB utilizes a novel \textit{Byte-Offset Anchor} architecture to enable zero-redundancy context expansion directly from compressed sidecars. This approach allows for efficient, sub-millisecond retrieval-augmented generation over multi-terabyte datasets without the memory overhead traditional vector databases require. Additionally, we aim to expand our structured extraction pipeline to include a broader range of SEC disclosures (e.g., DEF 14A, Form SR), providing a foundation for deep learning models capable of detecting material shifts in insider trading behavior. 

By leveraging deterministic XBRL parsers and LLM-agentic extraction, we can establish a programmatic Ground Truth that correlates insider transactions with specific vesting events and corporate liquidity cycles. This high-fidelity dataset will enable the development of sequence-based models (e.g., LSTMs or Temporal Fusion Transformers) capable of identifying subtle shifts in insider disclosure patterns that may signal material, non-public information. This architecture is also directly applicable to other high-stakes structured filings, such as Schedule 13D/13G and Form D, enabling the automated generation of massive, high-fidelity fine-tuning datasets across the global financial ecosystem.

\bibliographystyle{IEEEtran}
\bibliography{references}

\appendices
\section{Full Field-Level Extraction Metrics (Qwen3-Coder-30B-A3B)}
The following table provides the exhaustive precision, recall, and F1-score for all 137 fields evaluated in the 500-filing holdout set using the Qwen3-Coder-30B-A3B model.

\begin{small}
\begin{longtable}{lccc}
\caption{Exhaustive Field-Level Metrics (Qwen3-Coder-30B-A3B)} \label{tab:full_metrics} \\
\toprule
\textbf{Field Name} & \textbf{Precision} & \textbf{Recall} & \textbf{F1-Score} \\
\midrule
\endfirsthead
\toprule
\textbf{Field Name} & \textbf{Precision} & \textbf{Recall} & \textbf{F1-Score} \\
\midrule
\endhead
\bottomrule
\endfoot
derivative\_holdings\_conversion\_or\_exercise\_price & 0.453 & 0.930 & 0.609 \\
derivative\_holdings\_conversion\_or\_exercise\_price\_fn & 0.169 & 1.000 & 0.290 \\
derivative\_holdings\_direct\_or\_indirect\_ownership & 0.513 & 0.935 & 0.663 \\
derivative\_holdings\_direct\_or\_indirect\_ownership\_fn & 0.000 & 0.000 & 0.000 \\
derivative\_holdings\_exercise\_date & 0.966 & 1.000 & 0.983 \\
derivative\_holdings\_exercise\_date\_fn & 0.286 & 0.550 & 0.376 \\
derivative\_holdings\_expiration\_date & 0.778 & 1.000 & 0.875 \\
derivative\_holdings\_expiration\_date\_fn & 0.147 & 1.000 & 0.256 \\
derivative\_holdings\_nature\_of\_ownership & 0.872 & 1.000 & 0.932 \\
derivative\_holdings\_nature\_of\_ownership\_fn & 0.333 & 1.000 & 0.500 \\
derivative\_holdings\_security\_title & 0.513 & 1.000 & 0.678 \\
derivative\_holdings\_security\_title\_fn & 0.000 & 0.000 & 0.000 \\
derivative\_holdings\_shares\_owned\_following\_transaction & 0.487 & 1.000 & 0.655 \\
derivative\_holdings\_shares\_owned\_following\_transaction\_fn & 0.267 & 1.000 & 0.421 \\
derivative\_holdings\_transaction\_form\_type & 0.162 & 1.000 & 0.279 \\
derivative\_holdings\_underlying\_security\_shares & 0.453 & 0.930 & 0.609 \\
derivative\_holdings\_underlying\_security\_shares\_fn & 0.143 & 1.000 & 0.250 \\
derivative\_holdings\_underlying\_security\_title & 0.479 & 1.000 & 0.647 \\
derivative\_holdings\_underlying\_security\_value & 0.487 & 1.000 & 0.655 \\
derivative\_holdings\_value\_owned\_following\_transaction & 0.487 & 1.000 & 0.655 \\
derivative\_transactions\_conversion\_or\_exercise\_price & 0.982 & 0.944 & 0.963 \\
derivative\_transactions\_conversion\_or\_exercise\_price\_fn & 0.962 & 0.921 & 0.941 \\
derivative\_transactions\_deemed\_execution\_date & 1.000 & 0.986 & 0.993 \\
derivative\_transactions\_direct\_or\_indirect\_ownership & 0.996 & 0.938 & 0.966 \\
derivative\_transactions\_direct\_or\_indirect\_ownership\_fn & 0.077 & 0.500 & 0.133 \\
derivative\_transactions\_equity\_swap\_involved & 0.996 & 0.958 & 0.977 \\
derivative\_transactions\_exercise\_date & 0.957 & 0.971 & 0.964 \\
derivative\_transactions\_exercise\_date\_fn & 0.976 & 0.606 & 0.748 \\
derivative\_transactions\_expiration\_date & 0.958 & 0.903 & 0.929 \\
derivative\_transactions\_expiration\_date\_fn & 0.948 & 0.535 & 0.684 \\
derivative\_transactions\_nature\_of\_ownership & 0.895 & 0.959 & 0.926 \\
derivative\_transactions\_nature\_of\_ownership\_fn & 0.676 & 0.962 & 0.794 \\
derivative\_transactions\_price & 0.000 & 0.000 & 0.000 \\
derivative\_transactions\_security\_title & 0.986 & 0.944 & 0.965 \\
derivative\_transactions\_security\_title\_fn & 0.257 & 1.000 & 0.409 \\
derivative\_transactions\_shares & 0.000 & 0.000 & 0.000 \\
derivative\_transactions\_shares\_fn & 0.000 & 0.000 & 0.000 \\
derivative\_transactions\_shares\_owned\_following\_transaction & 0.975 & 0.937 & 0.956 \\
derivative\_transactions\_shares\_owned\_following\_transaction\_fn & 0.938 & 1.000 & 0.968 \\
derivative\_transactions\_transaction\_acquired\_disposed\_code & 0.982 & 0.941 & 0.961 \\
derivative\_transactions\_transaction\_code & 0.986 & 0.944 & 0.965 \\
derivative\_transactions\_transaction\_date & 0.993 & 0.951 & 0.972 \\
derivative\_transactions\_transaction\_date\_fn & 1.000 & 1.000 & 1.000 \\
derivative\_transactions\_transaction\_form\_type & 0.989 & 0.955 & 0.972 \\
derivative\_transactions\_transaction\_timeliness & 0.993 & 1.000 & 0.997 \\
derivative\_transactions\_transaction\_total\_value & 0.000 & 0.000 & 0.000 \\
derivative\_transactions\_underlying\_security\_shares & 0.989 & 0.641 & 0.778 \\
derivative\_transactions\_underlying\_security\_shares\_fn & 0.333 & 0.875 & 0.483 \\
derivative\_transactions\_underlying\_security\_title & 0.978 & 0.632 & 0.768 \\
derivative\_transactions\_underlying\_security\_title\_fn & 0.556 & 0.625 & 0.588 \\
derivative\_transactions\_underlying\_security\_value & 0.957 & 0.620 & 0.753 \\
derivative\_transactions\_underlying\_security\_value\_fn & 0.000 & 0.000 & 0.000 \\
derivative\_transactions\_value\_owned\_following\_transaction & 0.953 & 0.916 & 0.934 \\
footnotes\_1 & 1.000 & 0.986 & 0.993 \\
footnotes\_10 & 1.000 & 1.000 & 1.000 \\
footnotes\_11 & 1.000 & 1.000 & 1.000 \\
footnotes\_12 & 1.000 & 1.000 & 1.000 \\
footnotes\_2 & 0.990 & 0.976 & 0.983 \\
footnotes\_3 & 0.994 & 0.961 & 0.978 \\
footnotes\_4 & 0.989 & 0.959 & 0.974 \\
footnotes\_5 & 1.000 & 0.979 & 0.989 \\
footnotes\_6 & 1.000 & 0.963 & 0.981 \\
footnotes\_7 & 1.000 & 0.941 & 0.970 \\
footnotes\_8 & 1.000 & 1.000 & 1.000 \\
footnotes\_9 & 1.000 & 1.000 & 1.000 \\
form\_type & 1.000 & 0.988 & 0.994 \\
is\_rule\_10b5\_plan & 0.994 & 0.982 & 0.988 \\
issuer\_cik & 1.000 & 0.988 & 0.994 \\
issuer\_foreign\_trading\_symbol & 1.000 & 1.000 & 1.000 \\
issuer\_name & 1.000 & 0.988 & 0.994 \\
issuer\_schema\_version & 1.000 & 0.988 & 0.994 \\
issuer\_trading\_symbol & 0.996 & 0.988 & 0.992 \\
no\_securities\_owned & 1.000 & 0.988 & 0.994 \\
non\_derivative\_holdings\_direct\_or\_indirect\_ownership & 0.804 & 0.991 & 0.888 \\
non\_derivative\_holdings\_direct\_or\_indirect\_ownership\_fn & 0.429 & 0.600 & 0.500 \\
non\_derivative\_holdings\_nature\_of\_ownership & 0.899 & 0.938 & 0.918 \\
non\_derivative\_holdings\_nature\_of\_ownership\_fn & 0.783 & 0.878 & 0.828 \\
non\_derivative\_holdings\_security\_title & 0.776 & 0.996 & 0.872 \\
non\_derivative\_holdings\_security\_title\_fn & 0.222 & 1.000 & 0.364 \\
non\_derivative\_holdings\_shares\_owned\_following\_transaction & 0.762 & 0.995 & 0.863 \\
non\_derivative\_holdings\_shares\_owned\_following\_transaction\_fn & 0.440 & 1.000 & 0.611 \\
non\_derivative\_holdings\_transaction\_form\_type & 0.247 & 1.000 & 0.397 \\
non\_derivative\_holdings\_value\_owned\_following\_transaction & 0.762 & 0.995 & 0.863 \\
non\_derivative\_transactions\_deemed\_execution\_date & 1.000 & 1.000 & 1.000 \\
non\_derivative\_transactions\_direct\_or\_indirect\_ownership & 0.993 & 0.929 & 0.960 \\
non\_derivative\_transactions\_direct\_or\_indirect\_ownership\_fn & 0.438 & 0.824 & 0.571 \\
non\_derivative\_transactions\_equity\_swap\_involved & 1.000 & 0.989 & 0.994 \\
non\_derivative\_transactions\_nature\_of\_ownership & 0.940 & 0.950 & 0.945 \\
non\_derivative\_transactions\_nature\_of\_ownership\_fn & 0.982 & 0.982 & 0.982 \\
non\_derivative\_transactions\_price & 1.000 & 0.989 & 0.994 \\
non\_derivative\_transactions\_price\_fn & 0.920 & 0.983 & 0.950 \\
non\_derivative\_transactions\_security\_title & 1.000 & 0.989 & 0.994 \\
non\_derivative\_transactions\_security\_title\_fn & 1.000 & 1.000 & 1.000 \\
non\_derivative\_transactions\_shares & 1.000 & 0.989 & 0.994 \\
non\_derivative\_transactions\_shares\_fn & 0.848 & 0.979 & 0.909 \\
non\_derivative\_transactions\_shares\_owned\_following\_transaction & 0.990 & 0.979 & 0.985 \\
non\_derivative\_transactions\_shares\_owned\_following\_transaction\_fn & 0.979 & 0.969 & 0.974 \\
non\_derivative\_transactions\_transaction\_acquired\_disposed\_code & 1.000 & 0.989 & 0.994 \\
non\_derivative\_transactions\_transaction\_acquired\_disposed\_code\_fn & 0.833 & 1.000 & 0.909 \\
non\_derivative\_transactions\_transaction\_code & 1.000 & 0.989 & 0.994 \\
non\_derivative\_transactions\_transaction\_date & 1.000 & 0.989 & 0.994 \\
non\_derivative\_transactions\_transaction\_date\_fn & 0.857 & 1.000 & 0.923 \\
non\_derivative\_transactions\_transaction\_form\_type & 0.996 & 0.989 & 0.992 \\
non\_derivative\_transactions\_transaction\_timeliness & 0.996 & 1.000 & 0.998 \\
non\_derivative\_transactions\_transaction\_total\_value & 0.965 & 0.954 & 0.960 \\
non\_derivative\_transactions\_transaction\_total\_value\_fn & 0.000 & 0.000 & 0.000 \\
non\_derivative\_transactions\_value\_owned\_following\_transaction & 0.961 & 0.950 & 0.955 \\
non\_derivative\_transactions\_value\_owned\_following\_transaction\_fn & 0.000 & 0.000 & 0.000 \\
not\_subject\_to\_section\_16 & 1.000 & 0.988 & 0.994 \\
period\_of\_report & 1.000 & 0.988 & 0.994 \\
remarks & 0.998 & 0.998 & 0.998 \\
reporting\_owners\_ccc & 1.000 & 1.000 & 1.000 \\
reporting\_owners\_cik & 1.000 & 0.989 & 0.994 \\
reporting\_owners\_city & 0.986 & 0.977 & 0.982 \\
reporting\_owners\_country & 0.967 & 1.000 & 0.983 \\
reporting\_owners\_good\_address & 0.998 & 0.987 & 0.992 \\
reporting\_owners\_is\_10pctowner & 1.000 & 0.989 & 0.994 \\
reporting\_owners\_is\_director & 0.994 & 0.983 & 0.988 \\
reporting\_owners\_is\_director\_nominee & 1.000 & 0.989 & 0.994 \\
reporting\_owners\_is\_officer & 1.000 & 0.989 & 0.994 \\
reporting\_owners\_is\_other & 1.000 & 0.989 & 0.994 \\
reporting\_owners\_name & 0.998 & 0.987 & 0.992 \\
reporting\_owners\_non\_us\_address\_flag & 0.961 & 0.950 & 0.956 \\
reporting\_owners\_non\_us\_state\_territory & 0.966 & 1.000 & 0.982 \\
reporting\_owners\_officer\_title & 0.983 & 0.973 & 0.978 \\
reporting\_owners\_other\_text & 1.000 & 1.000 & 1.000 \\
reporting\_owners\_state & 0.992 & 0.985 & 0.988 \\
reporting\_owners\_state\_description & 0.500 & 0.923 & 0.649 \\
reporting\_owners\_zip & 0.992 & 0.985 & 0.988 \\
signatures\_signature\_date & 0.996 & 0.988 & 0.992 \\
signatures\_signature\_name & 0.994 & 0.987 & 0.990 \\
\end{longtable}
\end{small}

\end{document}